\begin{abstract}

Learned image priors are often proposed as replacements for hand-crafted regularizers in sparse-view computed tomography (CT), yet they are usually judged by PSNR on data resembling
their training set. We evaluate one noise-conditioned DnCNN denoiser used as a Plug-and-Play (PnP)
prior in proximal gradient against filtered back-projection (FBP), Tikhonov and total-variation
(TV) regularization, with validation-tuned parameters for every method and paired statistics,
under distribution shift, forward-model error and Poisson noise, and with a synthetic
lesion-insertion test (paired detection AUC, lesion retention, false lesion-like blobs).
On synthetic $28\times28$ digit images the prior beats a box-constrained TV by about $3.4$\,dB at
moderate noise ($\sigma=0.05$, 20 angles); the advantage shrinks at both ends of the noise range, the
ranking near the noise-free limit depends on how the data are simulated, and it reverses under distribution
shift: an MNIST-trained prior is $3.1$\,dB (Fashion-MNIST) and $4.3$--$4.6$\,dB (Shepp--Logan phantom)
below TV, whereas a Fashion-MNIST-trained prior beats TV by $2.2$\,dB on Fashion-MNIST, consistent with
prior mismatch rather than a failure of PnP itself. These results are reproduced under Poisson noise.
On digits the prior detects small lesions less reliably than TV even in distribution, and under
mismatch it creates lesion-like blobs more often than TV for most lesion types.
On real LoDoPaB-CT slices (patient-disjoint split; 30, 60 and 120 of 1000 angles; 4096 photons per
detector bin and one quarter and one tenth of that; 40 test patients) it does not consistently
outperform a well-tuned TV given the same iteration budget (300 iterations) in PSNR ($-0.93$ to $+0.33$\,dB over
nine conditions; at most $+0.49$\,dB when initialised from the TV solution) and has lower SSIM in eight of the
nine, but it detects small high-contrast lesions better than that TV in all nine (paired AUC $+0.03$ to $+0.16$,
six significant). It never significantly exceeds Hann-filtered FBP in detection, which the PSNR-tuned TV falls
below in all nine conditions (seven significant) despite a PSNR gain of $0.8$--$2.3$\,dB. On non-anatomical
phantoms its PSNR and SSIM deficits against TV are consistent and the detection advantage persists in one of
six conditions. PSNR and lesion detectability thus rank the methods differently, and a learned prior's value
depends on how well it matches the data it is applied to.
\end{abstract}
